\section{T4 Snapshot: Termination via well-founded rank (C1)}
\label{sec:t4-termination}

This note freezes the statement of the core ``outside-the-sandbox'' theorem \textbf{C1}:
\emph{termination from a strictly decreasing rank}.
The goal is a reusable, measure-theory-free kernel that supports many examples.

\paragraph{Lean source (as of T4).}
The canonical statement and its supporting lemmas live in
\texttt{Theorems/Termination.lean} (see also the extracted listing \texttt{Termination.txt}).

\subsection{Ranked objects (recap)}
Fix a type $X$ and a rank function $r : X \to \mathbb{N}$.
In Lean, this is a structure \texttt{ST.Ranked Nat X} with field \texttt{rank : X → Nat}.

\begin{definition}[Rank-descending relation]
\label{def:desc}
Given $R=(X,r)$, define the \emph{rank-descending} relation
\[
  \mathrm{Desc}_R(x,y) \;:\!\!\iff\; r(x) < r(y).
\]
\end{definition}

\begin{lemma}[Well-foundedness of $\mathrm{Desc}_R$]
\label{lem:wf-desc}
$\mathrm{Desc}_R$ is well-founded.
\end{lemma}

\noindent
The proof is the standard \emph{measure} argument: any relation induced by a function into $\mathbb{N}$
with strict $<$ is well-founded.

\subsection{The C1 assumption: strict rank decrease}
\begin{definition}[\texttt{RankDecreases}]
\label{def:rankdecreases}
Let $R=(X,r)$ and a step function $\mathrm{step}:X\to X$ be given.
We say \emph{\texttt{step strictly decreases rank}} if
\[
  \forall x\in X,\quad r(\mathrm{step}(x)) < r(x).
\]
In Lean this is the predicate \texttt{ST.RankDecreases R step}.
\end{definition}

\subsection{The C1 conclusion: termination / no infinite chain}
To express termination, we encode the ``one-step predecessor'' relation induced by \texttt{step}:

\begin{definition}[Step relation]
\label{def:steprel}
\[
  \mathrm{StepRel}_{\mathrm{step}}(x,y) \;:\!\!\iff\; x = \mathrm{step}(y).
\]
\end{definition}

\begin{theorem}[C1: well-founded step relation]
\label{thm:c1-wf-steprel}
If \texttt{RankDecreases R step} holds, then $\mathrm{StepRel}_{\mathrm{step}}$ is well-founded.
Equivalently, there is no infinite sequence
\[
  x_0,\; x_1,\; x_2,\;\dots
\]
such that $x_{n+1} = \mathrm{step}(x_n)$ for all $n$.
\end{theorem}

\paragraph{Proof idea (frozen).}
The well-foundedness proof is intentionally ``one-line'' and stable:
\[
  \mathrm{StepRel}(x,y)\Rightarrow x=\mathrm{step}(y)\Rightarrow r(x)<r(y).
\]
Thus $\mathrm{StepRel}$ is a subrelation of $\mathrm{Desc}_R$, and inherits well-foundedness.

\subsection{Convenience corollaries (non-goal, but useful)}
The Lean file also includes reusable corollaries:
\begin{itemize}
\item \texttt{wf\_of\_rank\_decreasing}: any relation whose steps strictly decrease rank is well-founded.
\item \texttt{no\_infinite\_desc}: a well-founded relation admits no infinite descending chain.
\item \texttt{acc\_of\_rank}: accessibility for $\mathrm{Desc}_R$ at every point.
\end{itemize}

\paragraph{What we intentionally do \emph{not} do at T4.}
We do not connect this theorem to the older ``Cat\_*'' families, nor do we generalize the codomain
$\mathbb{N}$ to an arbitrary well-founded order. Those belong to later tasks; C1 is frozen as the
minimal, practical kernel.
